\chapter{TWO WHEELED INVERTED PENDULUM IMPLEMENTATION}
\label{ch:platform}

\section{PHYSICAL DESIGN}

The robot consists of two mounting plates separated by threaded rods, with two
geared DC motors mounted on the lower plate and the electronics distributed
between the plates. The center of gravity sits above the wheel axis, giving an
inverted pendulum whose pivot is the wheel contact line. The physical parameters
are summarized in Table~\ref{tab:phys}.

\begin{figure}[htbp]
\centering
\includegraphics[width=0.45\textwidth]{twip_robot}
\caption{Two wheeled inverted pendulum robot.}
\label{fig:robot}
\end{figure}

\begin{table}[htbp]
\centering
\caption{TWIP robot physical parameters}
\label{tab:phys}
\begin{tabular}{lcc}
\toprule
Parameter & Value & Units \\
\midrule
Pendulum moment of inertia about the wheel axis, $I_p$ & 0.025\,$^a$ & \si{kg.m^2} \\
Wheel moment of inertia, $I_w$ & $1.441\times10^{-4}$\,$^b$ & \si{kg.m^2} \\
Distance to CG from wheel axis, $l$ & 0.075 & \si{m} \\
Mass of pendulum, $M_p$ & 1.432 & \si{kg} \\
Mass of wheels, $M_w$ & 0.1168 & \si{kg} \\
Wheel radius, $r$ & 0.045 & \si{m} \\
\bottomrule
\end{tabular}\\[2pt]
{\footnotesize $^a$ Computed with parallel-axis terms, i.e.\ about the wheel axis.
The equations of motion require the value about the center of gravity,
\SI{0.0169}{kg.m^2}; see Section~\ref{sec:simmodel}.\\
\rev{$^b$ The simulations, here and in the revision~7 programs, use the
uniform-disk value $M_wr^2/2=\SI{1.183e-4}{kg.m^2}$. The difference changes
$\beta$ of \eqref{eq:beta} by 1.4\%.}}
\end{table}

\rev{The moment of inertia and the CG offset are the two parameters least well
known on a build of this kind. Both are computed from a lumped-mass
idealization of the plates and components rather than measured. The gains that
balanced the hardware differ from the LQR gains computed from
Table~\ref{tab:phys}, but much of that difference is in the position and
velocity gains, whose sign the firmware's encoder frame affects
(Section~\ref{sec:gaintest}), so it does not measure the parameter error. The
recorded runs show the error nonetheless: at switch-on the tilt input gain is 8
to 26\% above the model (Section~\ref{sec:voltagetest}), and the accelerometer
points to an $I_p$ above the lumped-mass estimate
(Section~\ref{sec:momtest}). This is the motivation for the uncertainty
estimation problem rather than an incidental nuisance;
Appendix~\ref{app:validation} describes a pendulum-swing procedure for
measuring $I_p$ and $l$ directly.}

\section{ELECTRONIC DESIGN}

Electronics are a major component of any robot, and the two wheeled inverted
pendulum is no exception. The robot consists of a main computer, two motors,
various sensors, and power and communication devices.

\subsection{Computer}

The robot is controlled using an Arduino Due microcontroller board based on the
Atmel SAM3X8E ARM Cortex-M3, a 32-bit \SI{84}{MHz} CPU. The Due has
\SI{96}{kB} of onboard SRAM and \SI{512}{kB} of flash available for
programming. It provides 54 digital input/output channels, twelve of which can
be configured for pulse width modulation output, twelve analog inputs, and two
analog outputs, and supports serial communication over USB. The microcontroller
is programmed in C through the Arduino interface. The Arduino is a widely used
hobbyist platform and a broad range of integrated shields can be purchased to
work directly with it.

\rev{The Due is adequate for the control law but not generous. As documented in
Section~\ref{sec:difficulties}, serial transmission for telemetry dominated the
loop budget. Appendix~\ref{app:validation} recommends separating the
deterministic control path from the telemetry path so that the achievable loop
rate is set by the control computation rather than by the logging.}

\subsection{Motors}

Two Pololu \SI{12}{V} \SI{37}{mm} brushed DC motors drive the wheels. The
motors have a 30:1 gearbox and 64 counts-per-revolution integrated Hall effect
quadrature encoders measuring the position of the motor output shaft. The motors
are characterized by their free-run speed and current and their stall torque and
current; from these the torque constant $k_m$ and the back-EMF constant $k_e$
are obtained. The nominal terminal resistance is measured directly across the
output terminals at several shaft positions and averaged. The motor constants
are summarized in Table~\ref{tab:motor}.

Motor drive is provided by a Pololu Dual VNH5019 Motor Driver Arduino Shield.
The VNH5019 chip allows DC motor operation between \SI{5.5}{V} and \SI{24}{V} at
currents up to \SI{12}{A}, provides current sensing, and supports PWM operation
at ultrasonic frequencies up to \SI{20}{kHz} for quieter operation. The shield
is designed to work directly with the Arduino platform and a supplied library
provides straightforward control of two independent motors.

\begin{table}[htbp]
\centering
\caption{DC motor physical parameters}
\label{tab:motor}
\begin{tabular}{lcc}
\toprule
Parameter & Value & Units \\
\midrule
Free-run speed & 350 & RPM \\
Free-run current & 300 & mA \\
Stall torque & 110 & oz-in \\
Stall current & 5000 & mA \\
Torque motor constant, $k_m$ & 0.11541 & \si{N.m/A} \\
Back-EMF constant, $k_e$ & 0.307\,$^b$ & \si{V.s/rad} \\
Nominal resistance, $R$ & 2.5 & \si{\ohm} \\
\bottomrule
\end{tabular}\\[2pt]
{\footnotesize $^b$ From the free-run data above. Revision~7 used
\SI{0.00361}{V.s/rad}, which is inconsistent with those data by a factor of
about 85; see Section~\ref{sec:simmodel}.}
\end{table}

\subsection{Sensors}

Sensors are necessary in any robot in order to measure the state of the system,
and the state must be known in order to properly control it. In this system the
measurements are obtained from an integrated inertial measurement unit and from
Hall effect digital encoders integrated on each motor.

\subsubsection{Inertial measurement unit.}

Integrated System in Package (SiP) microchips are widely used in consumer
electronics owing to their small size and low production cost. The InvenSense
MPU-9150 is a low cost, low power SiP motion tracking device marketed for
equipment such as smartphones and tablets. It includes a three-axis gyroscope, a
three-axis accelerometer, a three-axis magnetometer, and an onboard Digital
Motion Processor (DMP) intended to offload algorithms from the main processor.
The SparkFun MPU-9150 breakout board is used here, which places the surface
mount device on a board with the communication pins available for direct
soldering. The DMP and the magnetometer were not used; only the direct
accelerometer and gyroscope readings are required to obtain the necessary
states.

The gyroscope has a 16-bit digital output with user-selectable ranges from
\SI{\pm250}{\degree/s} to \SI{\pm2000}{\degree/s}. The lowest range is selected
because it has the highest sensitivity, \SI{131}{LSB\per(\degree\per s)},
corresponding to a resolution of $1/131$ or approximately
\SI{0.00763}{\degree/s}. Noise is specified by two parameters: a total RMS noise
of \SI{0.06}{\degree/s} at a \SI{92}{Hz} rate, and a noise spectral density of
\SI{0.005}{\degree/s/\sqrt{Hz}} at \SI{10}{Hz}. These are characterized more
carefully by direct test in Section~\ref{sec:noise}. The specifications are
summarized in Table~\ref{tab:gyro}.

The accelerometer also has a 16-bit digital output, reporting in multiples of
gravity. Calibration was required to obtain a \SI{1}{g} reading at rest. At the
minimum range of \SI{\pm2}{g} the sensitivity is \SI{16384}{LSB\per g}, an exact
resolution of $1/16384$ or approximately \SI{61.0}{\micro g}. Noise is specified
as \SI{4}{mg} RMS at a \SI{100}{Hz} rate and a power spectral density of
\SI{400}{\micro g/\sqrt{Hz}} at \SI{10}{Hz}. The difference in the reference
rates between the two sensors necessitates a direct analysis of the noise from
the units actually used, so that the simulation is representative. The
accelerometer specifications are summarized in Table~\ref{tab:accel}.

\begin{table}[htbp]
\centering
\caption{Gyroscope specifications}
\label{tab:gyro}
\begin{tabular}{lcc}
\toprule
Parameter & Value & Units \\
\midrule
Range & $\pm250$ & \si{\degree/s} \\
Sensitivity & 131 & LSB/(\si{\degree/s}) \\
Resolution & 0.00763 & \si{\degree/s} \\
Total RMS noise & 0.06 & \si{\degree/s} RMS \\
Power spectral density & 0.005 & \si{\degree/s/\sqrt{Hz}} \\
\bottomrule
\end{tabular}
\end{table}

\begin{table}[htbp]
\centering
\caption{Accelerometer specifications}
\label{tab:accel}
\begin{tabular}{lcc}
\toprule
Parameter & Value & Units \\
\midrule
Range & $\pm2$ & \si{g} \\
Sensitivity & 16384 & LSB/\si{g} \\
Resolution & 61.0 & \si{\micro g} \\
Total RMS noise & 4 & \si{mg} RMS \\
Power spectral density & 400 & \si{\micro g/\sqrt{Hz}} \\
\bottomrule
\end{tabular}
\end{table}

\subsubsection{Digital encoders.}

Each Pololu metal DC motor has an integrated Hall effect quadrature encoder
measuring the relative position of the motor output shaft. The shaft carries a
magnetic disk with 16 magnetic divisions, and two Hall effect sensors detect the
changing field as the shaft rotates. Counting both the rising and falling edges
of both signals gives a total resolution of 64 counts per revolution of the
motor output shaft. \rev{The two signals are a quarter cycle (\SI{90}{\degree})
apart, in quadrature, which is the origin of the term and what allows the
direction of rotation to be decoded.} In
combination with the 30:1 gearbox, the gearbox output has a measurable
resolution of 1920 counts per revolution. With \SI{90}{mm} diameter wheels the
encoder measures discrete steps of \SI{1.47e-4}{m}.

\rev{The encoders are mounted on the motor shaft, ahead of the gearbox. Gear
backlash is therefore invisible to the encoder: the measured shaft angle
advances while the wheel does not. This is a structural measurement limitation,
not a resolution limitation, and it cannot be removed by increasing the sample
rate. Appendix~\ref{app:validation} discusses load-side encoding as the remedy.}

\subsection{Miscellaneous Electronics}

Four potentiometers are connected to the Arduino Due to allow easy adjustment of
parameters, in this case the controller gains, and a switch is used to enable
and disable the controller. Miscellaneous level-shifting electronics are
included to permit communication with the Due, as is an HC-05 Bluetooth module
for wireless transmission of serial output data. The robot is powered by a
single high-discharge \SI{11.1}{V} 20C lithium polymer battery.

\section{PROGRAMMING}

On initialization the various pins and connections are configured, enabling
communication between the sensors, switches, Bluetooth module, and the
microcontroller. The motor controller is initialized and the IMU is enabled. On
initialization of the IMU the constant bias is removed from the sensor
measurements by averaging a run of samples taken with the robot stationary.

\subsection{Measurements}
\label{sec:measurements}

The accelerometer measures the specific force vector, which at rest is the
gravitational field vector expressed in body axes. Writing the measured vector
as $\mathbf{G}_p=[a_x\ a_y\ a_z]^T$ and relating it to the gravitational field
through a rotation sequence in roll $\varphi$ and pitch $\vartheta_p$
\cite{kuipers1999quaternions},
\begin{equation}\label{eq:gravvec}
\frac{\mathbf{G}_p}{\norm{\mathbf{G}_p}}
=\frac{1}{\sqrt{a_x^2+a_y^2+a_z^2}}
\begin{bmatrix}a_x\\a_y\\a_z\end{bmatrix}
=\begin{bmatrix}
-\sin\vartheta_p\\
\cos\vartheta_p\sin\varphi\\
\cos\vartheta_p\cos\varphi
\end{bmatrix},
\end{equation}
which can be solved for the roll and pitch angles,
\begin{equation}\label{eq:rollpitch}
\varphi=\tan^{-1}\!\left(\frac{a_y}{a_z}\right),
\qquad
\vartheta_p=\tan^{-1}\!\left(\frac{-a_x}{\sqrt{a_y^2+a_z^2}}\right).
\end{equation}

The relationship in \eqref{eq:rollpitch} is only one way to compute the pitch
from the accelerometer. Rotation matrices do not commute, but as only the pitch
angle is of concern a different rotation sequence may be used to obtain the same
pitch measurement. Computation on the microcontroller must be minimized wherever
possible, so by using the rotation sequence $R_{yxz}$ in the same manner an
easier relationship for the tilt angle is obtained,
\begin{equation}\label{eq:tiltaccel}
\theta_a=\tan^{-1}\!\left(\frac{a_x}{a_z}\right).
\end{equation}

This calculation assumes two things: that the initial orientation of the
accelerometer is aligned with the gravitational field vector along the $z$ axis,
and that there are no external accelerations. The second clearly fails when the
robot is moving, so filtering of this measurement is necessary. A complementary
filter \cite{higgins1975comparison} is used. The same task could be performed
by a Kalman filter, but the computational limitations of the microcontroller
make the complementary filter the appropriate choice. The filter combines the
accelerometer and gyroscope measurements to obtain a more accurate angle
estimate using a simplified high-pass and low-pass pair. The form implemented on
the robot is
\begin{equation}\label{eq:compfilter}
\theta_{k+1}=\lambda_c\bigl(\theta_k+g_y\Delta t\bigr)+\bigl(1-\lambda_c\bigr)\theta_a ,
\end{equation}
where $\lambda_c$ is a design parameter set to $0.99$, $g_y$ is the $y$-axis
gyroscope reading, $\theta_a$ is the accelerometer tilt angle from
\eqref{eq:tiltaccel}, and $\theta_k$ is the filter angle estimate. The
complementary filter is simple to implement and well suited to an
implementation where computation is limited.

The tilt angle is the only measurement that must be pre-filtered before the
DMSO or the Kalman filter are used to estimate the full state. The gyroscope
measurements are used directly for the tilt rate. The motor encoder outputs are
processed using hardware interrupts: each time a Hall effect sensor signal
changes state the interrupt fires and increments or decrements the encoder
position. Velocity is obtained by differencing the position readings.

\rev{\begin{remark}[The pre-filter is part of the measurement model]
Equation \eqref{eq:compfilter} is a dynamic element placed between the sensors
and the estimator, so the quantity presented to the DMSO as $y_k$ is not a raw
measurement but a filtered one with its own lag and its own error
characteristics. Two consequences follow. First, the measurement noise
covariance used by the Kalman filter baseline must describe the complementary
filter output, not the accelerometer. Second, the comparison reported in
Section~\ref{sec:implresults}, in which the complementary filter outperforms
both model-based estimators fed raw accelerometer tilt, is a statement about
this pre-filtering step rather than about the estimators.
\end{remark}}

\section{DIFFICULTIES}
\label{sec:difficulties}

During assembly and programming many difficulties were encountered. Initially
the motors were a substantial problem and were eventually replaced. The original
motors were \SI{6}{V} DC units with lower torque output and an enormous deadzone
nonlinearity, giving a correspondingly small range of controllable output.
Simulation studies showed that the best achievable behavior with those motors
was a large stable oscillation, and sensor noise made even that unachievable.
It was during this period that the potentiometers were added, so that controller
gains could be adjusted quickly without reprogramming the microcontroller.

The motors were replaced with larger \SI{12}{V} units that do not have the same
large deadzone and allow a wider output range for control. This created a
problem with the plastic mounting plates, which twisted under the additional
weight until stiffening rods were added to the underside. It also exposed an
initial design flaw: plastic rods were used to separate the two mounting plates,
and the resulting flexure allowed vibrations to contaminate the sensor
measurements. Steel threaded rods were substituted and stiffened the structure
satisfactorily.

The robot is limited in available computation, so computations were minimized
wherever possible. Serial output proved to be one of the largest consumers of
computational effort. Serial data is essential for evaluating performance, so
its content had to be selected carefully and limited in order to keep the loop
time acceptable. The minimum loop execution time obtained was \SI{0.01}{s}.

The LQR gains had to be modified from the values computed from the measured
parameters of the robot, which \rev{revision~7 took as evidence of}
inaccurately measured parameters. The model is also not entirely accurate to the true system: the
pendulum is modeled as a point mass attached by a thin beam to the motor axis,
whereas the mounting plates create more complex interactions than this.
\rev{Once the gains were retuned the robot balanced, for up to \SI{21}{s},
although none of the recorded gain sets is linearly stable on any candidate
model (Table~\ref{tab:gaintest}).}
\rev{The logs point to a more specific cause (Section~\ref{sec:gaintest}). The
switch-on transients indicate that the firmware's encoder position is mirrored
relative to the model. In that frame, position and velocity gains entered with
the signs of the LQR design act with the opposite sign on the robot, and the
position and velocity potentiometers, which the firmware maps to $[-15,0]$,
could not set the signs the design requires. Both v9 runs balanced on the tilt
gains, with the base held within a few millimeters; in the two v8 runs the base
drove away until the robot was switched off (test~1) or until a chatter ended
the run (test~2). In closed-loop simulation the corrected model reproduces the
v9 runs in that frame only if the floor resisted rolling by a few percent of
the weight, so the frame is not yet confirmed. Either way, the modification of
the gains is not evidence of parameter error alone.}

\rev{\begin{remark}[These are the case for the work]
Three of the observations above are the empirical case for
Chapters~\ref{ch:dmso} and~\ref{ch:control} rather than incidental setbacks. The
parameters are uncertain, which motivates online uncertainty estimation. The
actuator is nonlinear in a way the linear model does not capture, which
motivates the unmatched uncertainty formulation. The available computation
constrains the achievable sample rate, which is an implementation constraint
treated as a design variable in Appendix~\ref{app:validation} rather than as a
fixed property of the platform.
\end{remark}}
